\section{Nhiệm vụ 1. Dự đoán quan hệ trong đồ thị cảnh}

\subsection{Bài toán và đầu vào, đầu ra}
Trong Hình~\ref{fig:motifs}, một phụ nữ và một chiếc xe máy đã được khoanh vùng.
Biết hai đối tượng này có mặt trong ảnh vẫn chưa đủ để mô tả cảnh. Quan hệ
\textit{riding} (đang cưỡi) cho biết người phụ nữ đang làm gì với chiếc xe.
Trong đồ thị cảnh, hai đối tượng là hai đỉnh, còn quan hệ này là một cạnh có
hướng từ người phụ nữ đến xe máy.

\begin{figure}[!ht]
\centering
\includegraphics[width=.83\columnwidth]{neural-motifs-fig1.pdf}
\caption{Minh họa đầu vào và đích dự đoán. Phía trên là ảnh và các hộp đối tượng.
Phía dưới là đồ thị chú thích, trong đó có bộ ba
\textit{woman--riding--motorcycle}. Trích Hình~1, trang~1 của
Zellers và cộng sự~\cite{zellers2018neuralmotifs}.}
\label{fig:motifs}
\end{figure}

Nhóm xét bài toán phân lớp quan hệ (\textit{Predicate Classification}, PredCls).
\textbf{Đầu vào} gồm ảnh RGB, hộp bao và nhãn lớp thật của từng đối tượng.
\textbf{Đầu ra} là danh sách bộ ba (đối tượng chủ thể, quan hệ, đối tượng đích),
kèm điểm tin cậy. Hai đối tượng được xác định theo từng hộp cụ thể. Một ảnh có
$N$ đối tượng tạo ra $N(N-1)$ cặp có hướng cần xét. Mô hình tự xác định cặp
nào có quan hệ và dự đoán nhãn cho cặp đó. Các cạnh có nhãn trong đồ thị
chú thích chỉ được dùng để huấn luyện và chấm kết quả.

\subsection{Dữ liệu và cách đánh giá}
VG150 là tập con của Visual Genome~\cite{krishna2017visualgenome}, gồm
150 lớp đối tượng và 50 lớp quan hệ. Nhóm chia dữ liệu theo ảnh. Sau khi
đối chiếu metadata và giữ những ảnh có hộp cùng chú thích quan hệ, nhóm
lấy 5.000 ảnh hợp lệ đầu tiên theo thứ tự lưu trong phần train làm
validation, giữ 57.723 ảnh còn lại để huấn luyện và giữ 26.446 ảnh hợp lệ
từ phần test để đánh giá. Các cặp đối tượng của một ảnh luôn thuộc cùng một tập.

Mỗi cặp có 51 điểm đầu ra, gồm 50 quan hệ và lớp nền. Khi đánh giá với
\textit{graph constraint} (GC), nhóm giữ quan hệ có xác suất cao nhất trong
50 quan hệ của mỗi cặp, sau đó xếp hạng tất cả các cặp trong cùng ảnh.
Một dự đoán khớp chú thích khi đúng nhãn quan hệ, đúng hai lớp đối tượng theo
thứ tự và độ giao trên hợp (IoU) của từng hộp với hộp chuẩn đạt ít nhất 0,5.

\textbf{R@50} đo tỷ lệ quan hệ chuẩn được tìm lại trong 50 dự đoán đứng đầu
của mỗi ảnh, rồi lấy trung bình qua ảnh. \textbf{mR@50} tính recall theo ảnh
cho từng loại quan hệ có mặt, rồi lấy trung bình trên 50 lớp. Vì vậy, một
lượt tìm lại quan hệ hiếm có thể ít ảnh hưởng đến R@50 nhưng có ý nghĩa khi
xét mR@50. Toàn bộ bảng định lượng dùng cùng bộ chấm \textit{reference-v1}
trên test. Nhóm chấm lại các checkpoint đã chọn bằng validation. Tập test
này đã từng được đánh giá bằng bộ chấm trước đó.

\subsection{Hai mô hình cơ sở}
Mỗi đỉnh có vector 549 chiều. Vector gồm 512 chiều thị giác trích từ vùng
ảnh bằng CLIP ViT-B/32 đóng băng~\cite{radford2021clip}, năm số hình học
(tọa độ tâm, chiều rộng, chiều cao, diện tích đã chuẩn hóa) và 32 chiều
biểu diễn nhãn đối tượng học trong khi huấn luyện. Hai mô hình ở
Hình~\ref{fig:vg-arch} dùng cùng đầu vào và chấm cùng tập cặp.

\begin{figure}[!ht]
\centering
\begin{tikzpicture}[>=Latex,font=\footnotesize,align=center,
 box/.style={draw=inkblue,rounded corners=2pt,inner sep=4pt},
 model/.style={box,text width=.39\columnwidth,minimum height=12mm}]
\node[box] (x) {Đặc trưng từng đối tượng (549 chiều)};
\node[model,below left=5mm and 1mm of x.south] (m) {MLP\\Ghép chủ thể và đích};
\node[model,below right=5mm and 1mm of x.south] (g) {GAT\\Truyền tin trên đồ thị,\\rồi ghép chủ thể và đích};
\node[box,below=24mm of x.south] (y) {51 điểm cho mỗi cặp có hướng};
\draw[->] (x.south) -- (m.north);
\draw[->] (x.south) -- (g.north);
\draw[->] (m.south) -- (y.north);
\draw[->] (g.south) -- (y.north);
\end{tikzpicture}
\caption{Hai đường dự đoán quan hệ. Thứ tự chủ thể và đối tượng đích được
giữ khi ghép vector. Với GAT, đặc trưng của mỗi đỉnh còn phụ thuộc các đỉnh
khác trong cùng ảnh.}
\label{fig:vg-arch}
\end{figure}

\textbf{MLP.} Ghép hai vector theo thứ tự chủ thể và đích thành vector 1.098
chiều. Hai lớp tuyến tính có kích thước $1.098\to256\to51$, với ReLU ở giữa,
tạo điểm quan hệ chỉ từ cặp đang xét.

\textbf{GAT.} Mô hình chiếu mỗi đỉnh từ 549 xuống 256 chiều, sau đó dùng
hai lớp graph attention~\cite{velickovic2018gat}, mỗi lớp bốn đầu chú ý.
Cạnh truyền tin nối mọi cặp đối tượng khác nhau trong ảnh và có vòng tự nối (\textit{self-loop}).
Tại mỗi lớp, đỉnh tổng hợp thông tin từ các đỉnh lân cận theo trọng số chú ý.
Hai vector sau truyền tin được ghép thành 512 chiều và đi qua bộ phân lớp
$512\to256\to51$. Mô hình dùng dropout 0,1.

\textbf{Huấn luyện.} Nhóm dùng AdamW với tốc độ học $10^{-3}$, weight decay
$10^{-2}$, batch hiệu dụng 512 ảnh, tối đa 200 epoch và dừng sớm với
patience 15. Mỗi cặp dương được lấy kèm tối đa ba cặp nền. Vì hàm mất mát
nhận một nhãn đích, bộ lấy mẫu bỏ qua các cặp có nhiều nhãn quan hệ khác nhau
khi train. Bộ chấm vẫn giữ các chú thích này. Nhóm đánh giá 54 checkpoint, tương ứng 18 cấu hình, mỗi cấu hình chạy với ba seed
42, 123 và 3407.

\subsection{Điều chỉnh hàm mất mát và bổ sung tri thức nền}
\textbf{Điều chỉnh theo tần suất lớp.} Mười loại quan hệ phổ biến nhất chiếm
88,40\% chú thích quan hệ trong tập train. Để giảm thiên lệch về những lớp
này, nhóm so sánh cross-entropy (CE) với \textit{logit-adjusted loss} theo công thức
\[
\mathcal{L}_{\mathrm{adj}}(z,y)=\operatorname{CE}(z+0{,}25\log\pi,\,y).
\]
Ở đây $z$ là vector điểm trước softmax, $y$ là nhãn đích và $\pi$ là phân bố
lớp ước lượng từ train theo quy tắc lấy mẫu, có cộng một vào số đếm mỗi lớp.
Vì $\log\pi$ âm hơn với lớp hiếm, mô hình phải học điểm $z$ cao hơn cho lớp
đó để dự đoán đúng khi train. Khi suy luận, nhóm xếp hạng từ điểm $z$ gốc.

\textbf{Tri thức tại đỉnh.} Nhóm liên kết nhãn đối tượng với khái niệm trong
WordNet hoặc ConceptNet~\cite{miller1995wordnet,speer2017conceptnet}. Với nhãn
có liên kết, nhóm ghép một vector tri thức 64 chiều đã đóng băng và một cờ
hợp lệ vào đặc trưng đỉnh. WordNet liên kết được 150/150 lớp, ConceptNet
liên kết được 138/150 lớp. Nhóm dùng TransE cho cả hai nguồn, đồng thời thử
RotatE và ComplEx trên ConceptNet~\cite{bordes2013transe,sun2019rotate,trouillon2016complex}.

\textbf{Tri thức tại cặp.} Với \textit{pair fusion}, hai vector tri thức của
chủ thể và đích tạo ra phần hiệu chỉnh cộng vào điểm quan hệ của GAT.
Một cổng học được điều khiển biên độ hiệu chỉnh. Với \textit{visual fusion},
cổng còn nhận đặc trưng GAT của hai đối tượng trong ảnh và điều tiết riêng
cho từng quan hệ. Để kiểm tra vai trò của ánh xạ lớp--khái niệm, nhóm giữ
cùng kiến trúc và tráo vector giữa 138 lớp ConceptNet có liên kết hợp lệ
(\textit{shuffled-KG}). Cả ba seed huấn luyện dùng cùng một hoán vị cố định.

\subsection{Kết quả chính và phân tích ablation}
\begin{table}[!ht]
\centering\small
\caption{Dự đoán quan hệ trên 26.446 ảnh test VG150. Điểm là phần trăm,
trung bình $\pm$ độ lệch chuẩn mẫu của ba seed. Mỗi nhóm có đối chứng riêng.}
\label{tab:vgscore}
\begin{tabular}{@{}lrr@{}}\toprule
Cấu hình & R@50 & mR@50\\\midrule
\multicolumn{3}{@{}l}{\textit{Mô hình cơ sở và tri thức tại đỉnh}}\\
MLP & $63{,}56\pm0{,}77$ & $16{,}95\pm1{,}52$\\
GAT & $63{,}87\pm0{,}73$ & $18{,}61\pm0{,}25$\\
GAT nhánh WordNet, tắt KG & $63{,}89\pm0{,}38$ & $18{,}44\pm0{,}29$\\
WordNet--TransE & $64{,}15\pm0{,}29$ & $17{,}84\pm0{,}09$\\
WordNet--TransE, bỏ cạnh & $64{,}11\pm1{,}15$ & $18{,}41\pm0{,}90$\\
ConceptNet--TransE & $64{,}10\pm0{,}27$ & $18{,}39\pm0{,}29$\\
ConceptNet--TransE, bỏ cạnh & $64{,}59\pm0{,}19$ & $18{,}15\pm0{,}28$\\
ConceptNet--RotatE & $64{,}21\pm0{,}37$ & $18{,}57\pm0{,}42$\\
ConceptNet--ComplEx & $64{,}37\pm0{,}18$ & $18{,}33\pm0{,}24$\\\midrule
\multicolumn{3}{@{}l}{\textit{Hàm mất mát $\times$ tri thức tại cặp}}\\
CE, tắt KG & $64{,}42\pm0{,}10$ & $18{,}11\pm0{,}16$\\
CE, bật KG & $63{,}63\pm0{,}14$ & $19{,}27\pm0{,}75$\\
Adjusted, tắt KG & $62{,}18\pm0{,}16$ & $24{,}59\pm0{,}93$\\
Adjusted, bật KG & $62{,}01\pm0{,}31$ & $24{,}67\pm1{,}07$\\\midrule
\multicolumn{3}{@{}l}{\textit{Fusion với adjusted loss}}\\
GAT, tắt nhánh hiệu chỉnh & $62{,}06\pm0{,}39$ & $24{,}32\pm0{,}91$\\
Pair, KG thật & $62{,}37\pm0{,}19$ & $\mathbf{25{,}24\pm0{,}39}$\\
Pair, KG xáo trộn & $62{,}34\pm0{,}18$ & $25{,}00\pm0{,}71$\\
Visual, KG thật & $62{,}34\pm0{,}21$ & $24{,}52\pm0{,}64$\\
Visual, KG xáo trộn & $62{,}20\pm0{,}23$ & $24{,}82\pm1{,}25$\\
\bottomrule\end{tabular}
\end{table}

Trên cùng đặc trưng CLIP, GAT tăng mR@50 so với MLP ở cả ba seed, trung bình
1,65 điểm phần trăm. R@50 tăng 0,31 điểm (Bảng~\ref{tab:vgscore}).
Để xác định thành phần nào tạo nên chênh lệch, nhóm lần lượt thay một yếu tố
trong từng họ mô hình. Bảng~\ref{tab:vgdelta} ghi hiệu số ghép cặp theo seed,
tính trước khi làm tròn.

\begin{table}[!ht]
\centering\small\setlength{\tabcolsep}{3pt}
\caption{Ablation trên VG150 test. Dấu dương biểu thị cấu hình bên trái
cao hơn đối chứng bên phải. Các độ chênh tính bằng điểm phần trăm.}
\label{tab:vgdelta}
\begin{tabularx}{\columnwidth}{@{}Xrr@{}}\toprule
Phép so sánh & $\Delta$R@50 & $\Delta$mR@50\\\midrule
WordNet tại đỉnh: bật KG trừ tắt KG & $+0{,}26$ & $-0{,}60$\\
WordNet: giữ cạnh trừ bỏ cạnh & $+0{,}04$ & $-0{,}57$\\
ConceptNet: giữ cạnh trừ bỏ cạnh & $-0{,}49$ & $+0{,}24$\\
Adjusted trừ CE, tắt KG tại cặp & $-2{,}24$ & $+6{,}48$\\
Adjusted trừ CE, bật KG tại cặp & $-1{,}61$ & $+5{,}40$\\
Bật trừ tắt KG tại cặp, adjusted loss & $-0{,}17$ & $+0{,}09$\\
Pair: KG thật trừ tắt hiệu chỉnh & $+0{,}31$ & $+0{,}93$\\
Pair: KG thật trừ KG xáo trộn & $+0{,}02$ & $+0{,}24$\\
Visual: KG thật trừ KG xáo trộn & $+0{,}14$ & $-0{,}30$\\
\bottomrule\end{tabularx}
\end{table}

\textit{(1) Ghép KG tại đỉnh.} Khi giữ cấu hình WordNet--TransE và bật vector
tri thức, mR@50 giảm ở cả ba seed. Như vậy độ phủ liên kết 150/150 lớp vẫn
chưa đủ để cải thiện khả năng nhận biết quan hệ theo lớp.

\textit{(2) Truyền tin giữa các đối tượng.} Thí nghiệm \textit{edges-off} bỏ
cạnh giữa các đỉnh, giữ self-loop và giữ nguyên toàn bộ cặp cần chấm điểm.
Giữ cạnh giúp ConceptNet--TransE tăng 0,24 điểm mR@50, trong khi WordNet--TransE
giảm 0,57 điểm. Tác dụng của ngữ cảnh đồ thị phụ thuộc cách biểu diễn đỉnh.

\textit{(3) Loss và KG.} Trong thiết kế $2\times2$, thay CE bằng adjusted
loss tăng mR@50 cả khi tắt lẫn khi bật KG tại cặp. Khi đã dùng adjusted
loss, bật thêm KG chỉ tăng trung bình 0,09 điểm và có một seed giảm.
Thay đổi lớn nhất trong nhóm đối chứng này đến từ điều chỉnh loss, đi kèm
đánh đổi về R@50.

\textit{(4) Nội dung tri thức trong fusion.} Pair fusion với KG thật đạt
mR@50 cao nhất trong 18 cấu hình. So với đối chứng tráo vector, lợi thế
chỉ còn 0,24 điểm, với hai seed tăng và một seed giảm. Visual fusion cho
chênh lệch theo chiều ngược lại. Với ba seed và một hoán vị, lợi ích riêng
của việc gắn đúng tri thức cho từng lớp còn chưa ổn định.

\subsection{Quan hệ hiếm và phân tích lỗi}
\textbf{Quan hệ nào được cải thiện?} Nhóm xếp 50 quan hệ theo số lần xuất hiện
trong train, rồi chia thành 10 lớp phổ biến, 20 lớp trung bình và 20 lớp hiếm.
Ở hai cấu hình cùng tắt KG, thay CE bằng adjusted loss làm mR@50 của nhóm
trung bình tăng từ 16,19 lên 26,56\%, nhóm hiếm tăng từ 3,54 lên 9,55\%,
còn nhóm phổ biến giảm từ 51,10 xuống 50,72\%. Kết quả này giải thích vì sao
recall trung bình theo lớp tăng rõ dù recall trung bình theo ảnh giảm.

\textbf{Bộ ba chưa gặp trong train.} Tập test chứa 7.848 chú thích trên
4.660 ảnh có tổ hợp (lớp chủ thể, quan hệ, lớp đích) chưa xuất hiện trong
chú thích train. Nhóm giữ cách xếp hạng toàn ảnh và tính recall riêng cho
các chú thích này (zR@50). Với cùng quy tắc GC, chín checkpoint thuộc ba cấu
hình fusion đạt trung bình 11,71\% khi tắt hiệu chỉnh, 11,43\% với KG thật
tại cặp và 11,26\% với KG xáo trộn. Ưu thế mR@50 của pair fusion trên toàn
tập vì thế chưa chuyển thành ưu thế trên các tổ hợp mới.

\begin{figure}[!ht]
\centering
\includegraphics[width=.48\columnwidth]{hit03.png}\hfill
\includegraphics[width=.48\columnwidth]{miss03.png}
\caption{Hai cặp trong cùng ảnh validation số 89, chấm bằng GAT seed 42.
Hộp đỏ đánh dấu chủ thể, hộp xanh đánh dấu đối tượng đích. Trái:
\textit{flower--in--vase}, hạng 1. Phải: \textit{shelf--on--table}, hạng 64.}
\label{fig:vg-paired}
\end{figure}

\textbf{Một ảnh, hai kết quả.} Ở Hình~\ref{fig:vg-paired}, GAT chọn đúng
nhãn cho cả hai cặp, nhưng chỉ bộ ba ở hạng 1 được tính trúng trong top-50.
Bộ ba ở hạng 64 được thu hồi khi tăng ngưỡng lên 100. Một cơ chế bỏ sót khác
xuất hiện ở ảnh số 3. Cặp \textit{girl--hair} có cả hai nhãn chuẩn \textit{has}
và \textit{with}. GC chỉ giữ \textit{with}. Hai ví dụ cho thấy cần xem cả
điểm xếp hạng và số nhãn chuẩn của cặp trước khi quy lỗi cho bộ phân lớp.
Phụ lục~\ref{app:vgcases} liệt kê đủ năm trường hợp trúng và năm trường hợp
bỏ sót đã chọn trên validation.
